% State of the Art section. Paste into your own thesis, or \input it.
%
% Replaces thesis/introduction/03_motivation.tex, which held this argument temporarily in the
% Introduction (see clinical_background/05_hemodynamics.tex's header comment) pending this chapter's
% existence, per docs_thesis/state_of_the_art_plan.md. Rewritten against
% docs_thesis/find_research.md Stage 3's full-text corrections to griffo2026wss -- the earlier draft
% overstated what that paper shows; the corrected version is below. \label{sec:motivation} was not
% referenced anywhere else in the thesis, so this section uses a fresh label instead.
%
% Scope check against 01_framing.tex and 05_predicts_disease.tex, so this section adds rather than
% repeats: framing.tex states the mechanism in three sentences and cites asakura1990flow and
% griffo2026wss's headline equivalence only; predicts_disease.tex covers Glagov remodeling under
% "Two limits". This section is the one place the full mechanism (bifurcation and curvature
% separately) and the fully-caveated Griffo result belong.
%
% Verification status: griffo2026wss is full-text read (docs_thesis/papers/). The rest --
% murasato2022bifurcation, kahe2020tortuosity, wang2024lmphenotypes -- are abstract-level
% [BIBLIOGRAPHY VERIFIED] per docs_thesis/literature.md Axis 5 and docs_thesis/find_research.md.

\section{Why shape is enough}
\label{sec:shape-enough}

If geometry sets wall shear stress and shear sets the state of the endothelium
(Section~\ref{sec:hemodynamics}), then arteries that share a shape should share a shear pattern, and
therefore a similar baseline exposure to the mechanism that starts atherosclerosis.
Two separate mechanical routes connect a tree's shape to shear, and one recent result asks whether
the shear itself even needs to be computed once the shape is known.

A bifurcation splits one stream of blood into two, and the split sets the shear independently of the
angle at which the vessels diverge.
Murasato et al.\ reviewed how a dividing flow behaves at a coronary bifurcation
\cite{murasato2022bifurcation}: shear falls on the lateral wall of each daughter vessel and rises on
the carinal side, the wedge between them, in every geometry they examined.
What the angle changes is how sharply that gradient falls off, which is why a wide bifurcation
exposes more lateral wall to low shear than a narrow one.
Curvature works through a different route.
Flow through a bend develops a secondary, helical component, and the shear it produces runs low on
the inner wall of the bend and high on the outer one \cite{kahe2020tortuosity}.
Two different mechanisms, and both were confirmed on real tissue in Section~\ref{sec:framing}: the
outer wall of a vessel leaving a bifurcation and the inner wall of a curved segment are exactly where
plaque was found in five postmortem coronary trees \cite{asakura1990flow}.

If a phenotype maps onto a shear pattern this consistently, arteries sorted by shape alone should
separate by shear.
The closest direct test of that claim covers one anatomical site.
Wang et al.\ simulated flow through 76 left main arteries sorted by shape into four phenotypes
\cite{wang2024lmphenotypes}.
Shear fell lowest around the branch points of the LAD, and lowest of all in the widest-angled
phenotype.
The result is exactly what the mechanism above predicts: a geometric phenotype maps onto a
characteristic pattern of low shear, not just onto a single feature value.
It is also, so far, the only test of it.
The study covers one bifurcation in 76 arteries, simulated rather than measured in patients, and it
says nothing about whether the pattern distinguishes people who go on to develop plaque from people
who do not.

None of this requires a flow simulation to use, and the strongest test of that claim comes from a
diseased, not a healthy, cohort.
Griffo et al.\ trained a network to reproduce computational-fluid-dynamics shear from lumen geometry
alone, using 1078 coronary arteries reconstructed from invasive angiography in 748 patients pooled
across three trials of coronary stenosis, median 58.5\% area stenosis \cite{griffo2026wss}.
The network and the simulation it was trained to reproduce agreed almost exactly on which lesion in a
patient's tree went on to cause a heart attack: an area under the curve of 0.63 for both,
statistically indistinguishable (DeLong $p=0.84$), and 0.68 against 0.66 for the ratio between a
lesion's shear and its parent vessel's ($p=0.37$).
The authors call the discrimination itself modest, attributing the ceiling to how many other factors
besides shear decide whether a plaque ruptures; what carries weight here is the equivalence, not the
magnitude.
Three qualifications keep that equivalence from being read as more than it is.
The task is not population risk prediction: it discriminates 80 lesions that later became the
culprit from 107 that did not, within patients who had all already had a heart attack, so it answers
which lesion ruptured, not who will have an event.
The cohort was already diseased and reconstructed from invasive angiography, not from the
lower-resolution CCTA lumens this thesis and the Copenhagen General Population Study actually
supply, so whether the equivalence survives that transfer is untested.
And the spatial agreement the authors report between network and simulation, a Dice coefficient of
0.88, is measured on the highest-shear regions of the vessel, not on the low-shear tail the
mechanism above rests on.
What the result licenses is narrower than a method: if a geometric phenotype is defined at all, its
characteristic shear pattern does not need to be simulated to be inferred, because the geometry
already carries it.
That is the premise this thesis works from --- measuring shape directly, at population scale, rather
than solving the flow field lesion by lesion --- not a method being reproduced at that scale.
